\tightvspace{-5pt}
\section{Conclusion and Future Work}

Across image classification, text regression, and multi-label text classification, in-distribution accuracy turns out to neither guarantee nor predict temporal robustness. What governs the rate of degradation is instead the strength of a model's inductive bias and whether it relies on a frozen pretrained encoder. A stronger bias extracts the most discriminative, period-specific features and leads in distribution, yet those same features are the first to lose their meaning as the data drifts, so the architectures that fit the training period most tightly are the ones that degrade fastest. Frozen pretrained encoders occupy a different regime, conceding in-distribution accuracy in return for steadier behaviour over time, and on \texttt{arXiv}, where the task is already solved by the pretrained representation and the bias buys no advantage, no family decays faster than the rest. The practical reading is that the model with the best held-out score at training time is often the least robust once deployed, so architecture selection should weigh the expected horizon between retraining cycles, not in-distribution accuracy alone.

\tightvspace{-5pt}
\subsection{Limitations}

\tightvspace{-5pt}
Our comparison covers the most common inductive biases, convolutional, recurrent, and attention-based, against simple baselines and frozen encoders, but many architectures remain untested and the study is best read as a first systematic pass rather than an exhaustive one. It spans three datasets across two modalities, and broadening it to further domains and other forms of distribution shift would show how widely the pattern holds. The text tasks are the most constrained, using a narrow label space, a five-point rating for \texttt{Amazon Reviews} and seven categories for \texttt{arXiv}, spanning only about a decade and a quarter of a century against a full century for \texttt{Yearbook}, and running on a stratified subsample rather than the complete corpora. Because too few examples fall within each time slice to train a text encoder from scratch, the text models also share a single frozen \texttt{RoBERTa} representation, so their absolute scores should be read in that light. We average over a small fixed seed set rather than conducting formal significance tests, and hyperparameter choices and seed variance may still affect fine-grained rankings; we therefore emphasize family-level patterns over individual placements. The families are also not capacity-matched. The small, medium, and large variants within each family expose scale effects, and on \texttt{Yearbook} the smallest \texttt{MLP}, \texttt{CNN}, and \texttt{ResNet} variants decay less than their larger siblings, so inductive bias and capacity remain partially confounded. The analysis is also descriptive rather than mechanistic. We measure how robustness varies with inductive bias and pretraining without isolating the precise cues responsible, and we evaluate inherent robustness under cumulative training rather than any adaptive policy, which leaves the question of retraining orchestration to systems such as Modyn \cite{boether_modyn_2025}.

\tightvspace{-5pt}
\subsection{Future Work}

\tightvspace{-5pt}
These limitations point to several extensions. The comparison can be widened to a broader range of inductive biases, to capacity-matched pairs that disentangle bias from scale, and to additional datasets, modalities, and drift mechanisms, and carried to longer temporal spans with the full corpora and richer label spaces, where language drift should be more pronounced and the gap between architectures wider. The drift matrices can also be paired with scheduled or drift-triggered retraining under explicit compute budgets, measuring not only how much a model decays but how cheaply that decay can be undone \cite{mahadevan_cost_aware_2024}. Most of all, moving from description to mechanism, by tying inductive bias and pretraining back to the specific features that drift, would turn the observed regularity into an account of why temporal robustness behaves as it does.

\tightvspace{-5pt}
\subsection{Availability}

\tightvspace{-5pt}
Code, experiment definitions, and paper-generation scripts are available at \href{https://github.com/learning-mechanisms/drift-happens/}{github.com/learning-mechanisms/drift-happens}; public run histories and matrix artifacts at \href{https://wandb.ai/drift-happens/drift-happens}{wandb.ai/drift-happens/drift-happens}. The \extendedpreprintlink{} and \resultssitelink{} retain the complete per-model drift-matrix galleries. Code is Apache-2.0; paper, figures, and documentation are CC BY 4.0 where we hold the rights, and external datasets and models keep their original licenses.
